\section{From Recurrence to Attention}

The previous two sections developed a precise picture of what recurrent models do well and where they struggle.
Recurrence is a natural architecture for temporal data because it processes a sequence through an evolving state.
That state functions as a running summary of the past.
But this same mechanism also creates a structural bottleneck.
If information from an early part of a long sequence is needed much later, then it must survive repeated state updates and repeated Jacobian products.
Even gated architectures make that dependence travel through time one step after another.

\medskip

Attention arose as a structural response to this problem.
Instead of requiring all relevant information to be compressed into one fixed-size state, the model is allowed to access past representations directly and selectively.
The central idea is content-based retrieval inside the network: a model forms a query describing what it currently needs, compares that query with a collection of candidate keys, and then aggregates the corresponding values according to learned relevance weights.
The result is not merely a different formula.
It is a different architecture for memory access.

\medskip

The purpose of this section is to make that transition mathematically explicit.
We will define attention as content-based weighted aggregation, derive its weighted-sum form from similarity scores, explain the encoder--decoder bottleneck in classical sequence-to-sequence models, prove a basic structural proposition about the output of attention, and work through a small alignment example.
The pedagogical point is that attention did not appear as an arbitrary design fashion.
It solved a precise architectural problem.

\subsection*{The encoder--decoder bottleneck}

Consider a classical sequence-to-sequence architecture for translation.
An encoder reads an input sequence $x_1,\dots,x_T$ and produces a final state $h_T$.
A decoder then produces an output sequence $y_1,\dots,y_S$ using that final state as its initial summary of the source.
At an abstract level one writes
\[
 h_t = F(h_{t-1},x_t),
 \qquad
 s_u = G(s_{u-1},y_{u-1},h_T),
\]
where $h_t$ is the encoder state and $s_u$ is the decoder state.
The entire source sentence is therefore compressed into the single vector $h_T$.

\medskip

This architecture can work well for short sequences, but it imposes a severe information bottleneck.
A fixed-dimensional vector must simultaneously retain all the details that might later matter for every output token.
When sequences become longer or dependencies become more specific, the decoder may need information that was present earlier in the source but is not adequately represented in the final summary.
In that sense, classical sequence-to-sequence recurrence asks one vector to do too much.

\begin{figure}[t]
\centering
\begin{tikzpicture}[>=Latex, node distance=0.9cm and 0.8cm]
\tikzstyle{tok}=[draw, rounded corners, minimum width=0.9cm, minimum height=0.55cm]
\tikzstyle{state}=[draw, circle, minimum size=0.62cm, inner sep=0pt]
\node[tok] (x1) at (0,0) {$x_1$};
\node[tok] (x2) at (1.4,0) {$x_2$};
\node[tok] (x3) at (2.8,0) {$\cdots$};
\node[tok] (x4) at (4.2,0) {$x_T$};
\node[state] (h1) at (0,1.3) {$h_1$};
\node[state] (h2) at (1.4,1.3) {$h_2$};
\node[state] (h3) at (2.8,1.3) {$\cdots$};
\node[state, thick] (hT) at (4.2,1.3) {$h_T$};
\node[state] (s1) at (5.9,1.3) {$s_1$};
\node[state] (s2) at (7.3,1.3) {$s_2$};
\node[state] (s3) at (8.7,1.3) {$\cdots$};
\node[tok] (y1) at (5.9,0) {$y_1$};
\node[tok] (y2) at (7.3,0) {$y_2$};
\node[tok] (y3) at (8.7,0) {$y_S$};
\draw[->] (x1) -- (h1);
\draw[->] (x2) -- (h2);
\draw[->] (x4) -- (hT);
\draw[->, thick] (h1) -- (h2);
\draw[->, thick] (h2) -- (h3);
\draw[->, thick] (h3) -- (hT);
\draw[->, thick, line width=1pt] (hT) -- node[above] {single summary vector} (s1);
\draw[->, thick] (s1) -- (s2);
\draw[->, thick] (s2) -- (s3);
\draw[->] (s1) -- (y1);
\draw[->] (s2) -- (y2);
\draw[->] (s3) -- (y3);
\node[align=center] at (4.95,2.15) {encoder--decoder bottleneck};
\end{tikzpicture}
\caption{Classical sequence-to-sequence recurrence compresses the entire source into one summary state before decoding. Attention replaces this single bottleneck with selective access to many encoder states.}
\label{fig:seq2seq-bottleneck}
\end{figure}

Attention loosens this bottleneck.
Instead of forcing the decoder to rely only on $h_T$, it allows the decoder at time $u$ to construct a context vector from the entire collection $h_1,\dots,h_T$.
The context depends on what the decoder currently needs, so different output positions may retrieve different parts of the source.

\subsection*{Core definitions}

\begin{definition}[Attention mechanism]
Let $q\in\mathbb{R}^{d_q}$ be a query, let $k_1,\dots,k_n\in\mathbb{R}^{d_q}$ be keys, and let $v_1,\dots,v_n\in\mathbb{R}^{d_v}$ be values.
An \emph{attention mechanism} is a map that assigns a relevance weight $\alpha_i(q)$ to each pair $(q,k_i)$ and then returns an aggregate of the values:
\[
\operatorname{Attn}(q,\{(k_i,v_i)\}_{i=1}^n)
= \sum_{i=1}^{n} \alpha_i(q) v_i.
\]
When the weights depend on content similarity between the query and the keys, the mechanism is called \emph{content-based attention}.
\end{definition}

\begin{definition}[Compatibility score and alignment weights]
A \emph{compatibility score} or \emph{similarity score} is a function
\[
 s:\mathbb{R}^{d_q}\times \mathbb{R}^{d_q}\to\mathbb{R}
\]
that measures how relevant a key is to a given query.
Given scores $s(q,k_i)$, the corresponding \emph{alignment weights} are commonly defined by the softmax normalization
\[
 \alpha_i(q)=\frac{\exp(s(q,k_i))}{\sum_{j=1}^{n}\exp(s(q,k_j))},
 \qquad i=1,\dots,n.
\]
These weights are nonnegative and sum to $1$.
\end{definition}

\begin{definition}[Context vector]
In encoder--decoder attention, the \emph{context vector} at decoder step $u$ is the weighted combination
\[
 c_u = \sum_{t=1}^{T} \alpha_{u,t} h_t,
\]
where $h_t$ are encoder states and $\alpha_{u,t}$ are alignment weights determined by the current decoder state and the source representations.
The context vector summarizes the source \emph{relative to the current decoding need}, not in one fixed way for the whole sentence.
\end{definition}

\subsection*{From similarity scores to weighted aggregation}

The weighted-sum form of attention is not arbitrary.
It emerges naturally from three requirements.
First, each candidate memory item should receive a relevance score based on its match with the current query.
Second, these scores should be normalized so that they become comparable across different numbers of candidates.
Third, the retrieved information should remain differentiable with respect to both the scores and the stored representations.

\medskip

Suppose the current decoder state or query is $q$, and suppose the encoder has produced candidate memories $h_1,\dots,h_T$.
A score function $s(q,h_t)$ assigns a real number to each source position.
To transform these scores into usable weights, one applies a normalization map.
The most common choice is the softmax:
\[
 e_t = s(q,h_t),
 \qquad
 \alpha_t = \frac{e^{e_t}}{\sum_{j=1}^{T} e^{e_j}}.
\]
Now $\alpha_t\ge 0$ and $\sum_t \alpha_t=1$.
The natural differentiable retrieval rule is then
\[
 c(q)=\sum_{t=1}^{T}\alpha_t h_t.
\]
Thus the weighted-sum form of attention is obtained by the composition
\[
 \text{query and memories}
 \xrightarrow{\text{score}}
 (e_1,\dots,e_T)
 \xrightarrow{\text{softmax}}
 (\alpha_1,\dots,\alpha_T)
 \xrightarrow{\text{weighted sum}}
 c.
\]

\medskip

This formula has an important interpretation.
The model does not retrieve one memory slot in a hard, discrete way.
Instead it forms a soft data-dependent average over the available values.
That soft selection is what makes attention trainable by gradient-based methods.
At the same time, when one weight becomes much larger than the others, the mechanism approximates nearly discrete retrieval.

\subsection*{A structural proposition about attention outputs}

The weighted-sum form immediately imposes geometric structure on what attention can produce.

\begin{proposition}[Attention outputs lie in the span and, with softmax weights, in the convex hull of the values]
Let $v_1,\dots,v_n\in\mathbb{R}^{d_v}$ and let
\[
 y = \sum_{i=1}^{n}\alpha_i v_i.
\]
Then $y\in \operatorname{span}\{v_1,\dots,v_n\}$.
If in addition $\alpha_i\ge 0$ for all $i$ and $\sum_{i=1}^{n}\alpha_i=1$, then $y$ lies in the convex hull of $\{v_1,\dots,v_n\}$.
In particular, standard softmax attention produces outputs that are convex combinations of the available values.
\end{proposition}

\begin{proof}
By definition, $y$ is a linear combination of the vectors $v_1,\dots,v_n$, so it belongs to their span.
If moreover the coefficients are nonnegative and sum to $1$, then $y$ is a convex combination of those vectors, which is exactly the definition of membership in the convex hull.
For softmax attention the weights satisfy these properties automatically, so the conclusion follows.
\end{proof}

\medskip

This proposition is elementary but conceptually important.
Attention does not create information from nowhere.
It reweights and recombines existing value vectors.
Its power comes from making the coefficients data-dependent.
The architecture therefore changes \emph{which} information can be accessed and \emph{when} it can be accessed, even though each context vector remains a structured combination of the stored values.

\begin{remark}
The proposition also explains why learned linear projections are often placed before and after attention.
If one wants the retrieved vector to live in a more expressive feature space, one can transform the stored values before aggregation or transform the resulting context afterward.
Attention itself supplies adaptive selection; linear maps around it expand the representational class.
\end{remark}

\subsection*{Why attention solves a precise recurrent problem}

The bottleneck in classical sequence-to-sequence recurrence has two aspects.
First, the source must be summarized in a single vector before the decoder begins to use it.
Second, the contribution of an early source token to a late target token must travel through many recurrent transitions.
Attention addresses both problems.
At decoder step $u$, one computes a fresh query from the current decoding state and retrieves a context vector from all source positions.
Thus source token $x_t$ can influence output token $y_u$ directly through the coefficient $\alpha_{u,t}$ rather than only indirectly through a long chain of recurrent state updates.

\medskip

This does not remove all sequential structure.
In an encoder--decoder model the decoder still evolves through time.
But it replaces one global compressed memory with dynamic access to a whole collection of source-side representations.
That is the decisive conceptual change.
Memory is no longer only something carried forward; it becomes something that can be queried.

\subsection*{Worked example: a short alignment problem}

Consider a toy translation example with source sentence
\[
 \text{the cat sleeps}
\]
and target sentence
\[
 \text{le chat dort}.
\]
Suppose the encoder has produced source representations $h_1,h_2,h_3$ for \emph{the}, \emph{cat}, and \emph{sleeps}.
At the decoder step for producing \emph{chat}, the current decoder state should look most strongly at the source word \emph{cat}.
A simplified set of scores might be
\[
 e=(0.2,\,2.4,\,0.1).
\]
The softmax weights are then approximately
\[
 \alpha\approx(0.091,\,0.827,\,0.082).
\]
Therefore the context vector is
\[
 c \approx 0.091 h_1 + 0.827 h_2 + 0.082 h_3.
\]
This means that the decoder forms a source summary specialized to the current target word.
For the next decoder step, when producing \emph{dort}, the largest alignment weight would typically shift toward the source word \emph{sleeps}.

\begin{example}[A tiny numerical attention computation]
Suppose the three source values are one-dimensional for simplicity:
\[
 v_1=0,
 \qquad
 v_2=1,
 \qquad
 v_3=2,
\]
and suppose the attention weights at one decoder step are
\[
 \alpha=(0.1,0.8,0.1).
\]
Then the retrieved context is
\[
 c = 0.1\cdot 0 + 0.8\cdot 1 + 0.1\cdot 2 = 1.0.
\]
If the weights instead become $(0.05,0.15,0.80)$, then
\[
 c = 0.05\cdot 0 + 0.15\cdot 1 + 0.80\cdot 2 = 1.75.
\]
The values have not changed; what changed is the content-dependent routing.
This is the basic mechanism by which attention can focus on different parts of a sequence at different times.
\end{example}

\begin{figure}[t]
\centering
\begin{tikzpicture}[>=Latex]
% matrix labels
\node at (1.0,4.35) {the};
\node at (2.2,4.35) {cat};
\node at (3.4,4.35) {sleeps};
\node[rotate=90] at (-0.35,3.2) {le};
\node[rotate=90] at (-0.35,2.0) {chat};
\node[rotate=90] at (-0.35,0.8) {dort};
% cells row 1
\fill[black!20] (0.4,2.8) rectangle (1.6,4.0);
\fill[black!8]  (1.6,2.8) rectangle (2.8,4.0);
\fill[black!6]  (2.8,2.8) rectangle (4.0,4.0);
% row 2
\fill[black!10] (0.4,1.6) rectangle (1.6,2.8);
\fill[black!70] (1.6,1.6) rectangle (2.8,2.8);
\fill[black!8]  (2.8,1.6) rectangle (4.0,2.8);
% row 3
\fill[black!5]  (0.4,0.4) rectangle (1.6,1.6);
\fill[black!12] (1.6,0.4) rectangle (2.8,1.6);
\fill[black!78] (2.8,0.4) rectangle (4.0,1.6);
% grid
\draw[step=1.2cm, black] (0.4,0.4) grid (4.0,4.0);
% numbers
\node at (1.0,3.4) {0.72};
\node at (2.2,3.4) {0.16};
\node at (3.4,3.4) {0.12};
\node[text=white] at (2.2,2.2) {0.82};
\node at (1.0,2.2) {0.10};
\node at (3.4,2.2) {0.08};
\node at (1.0,1.0) {0.06};
\node at (2.2,1.0) {0.14};
\node[text=white] at (3.4,1.0) {0.80};
\end{tikzpicture}
\caption{A toy alignment matrix for translating \emph{the cat sleeps} to \emph{le chat dort}. Darker cells indicate larger attention weights. Different target words retrieve different source positions, which is exactly what a single fixed source summary cannot do.}
\label{fig:alignment-matrix}
\end{figure}

The alignment matrix in \Cref{fig:alignment-matrix} visualizes the same idea at the sentence level.
Each row corresponds to one decoder step, and each column corresponds to one source position.
The weights in each row sum to $1$.
For \emph{chat}, the mass is concentrated on \emph{cat}; for \emph{dort}, it is concentrated on \emph{sleeps}.
Thus the decoder forms different context vectors for different output tokens.
This is what it means for attention to provide dynamic, content-based access.

\subsection*{Relation to the later transformer viewpoint}

Pedagogically it is important to pause here before moving to transformers.
Attention should first be understood as a mechanism:
\begin{itemize}
    \item it responds to the encoder--decoder bottleneck,
    \item it replaces one fixed summary by query-dependent retrieval,
    \item it creates direct interactions between distant positions,
    \item it remains differentiable because the retrieval is soft rather than hard.
\end{itemize}
Only after this mechanism is clear does it become natural to study architectures built primarily from attention, such as transformers.
In that sense, attention is not yet the whole architecture.
It is the key structural idea that makes a different architecture possible.

\whattoremember{Attention is content-based weighted aggregation.
A query is compared with a collection of keys, the resulting scores are normalized into alignment weights, and a context vector is formed as a weighted sum of values.
This directly addresses the classical sequence-to-sequence bottleneck by allowing different output positions to retrieve different source information instead of relying on one fixed summary vector.}

\commonconfusion{Attention is sometimes described as if it were a mysterious replacement for recurrence.
That is too vague.
Its original role in sequence-to-sequence models was much more specific: it allowed the decoder to access the full collection of encoder states instead of depending only on a single compressed vector.
Also, attention is not hard lookup.
With softmax normalization it produces a convex combination of values, so its power comes from adaptive weighting rather than from magically creating new information.}

\subsection*{Exercises}

\begin{exercise}
Let values $v_1,v_2,v_3\in\mathbb{R}^d$ and weights $\alpha_1,\alpha_2,\alpha_3$ satisfy $\alpha_i\ge 0$ and $\alpha_1+\alpha_2+\alpha_3=1$.
Show directly that $\sum_i \alpha_i v_i$ lies in the convex hull of $\{v_1,v_2,v_3\}$.
Why does this apply automatically to softmax attention?
\end{exercise}

\begin{exercise}
Suppose a decoder step produces similarity scores $(1,2,0)$ against three source positions.
Compute the corresponding softmax weights.
Which source position receives the largest share of attention?
\end{exercise}

\begin{exercise}
Explain in your own words why the vector $h_T$ in a classical encoder--decoder model is a bottleneck.
Your answer should refer both to representation capacity and to long-range dependence through recurrent updates.
\end{exercise}

\begin{exercise}
In a toy translation example with source words $(x_1,x_2,x_3)$ and target words $(y_1,y_2)$, suppose the alignment weights are
\[
(0.7,0.2,0.1)\quad\text{for }y_1,
\qquad
(0.1,0.1,0.8)\quad\text{for }y_2.
\]
Describe qualitatively what this says about the different context vectors used at the two decoder steps.
\end{exercise}

\begin{exercise}
Why is the weighted-sum form of attention more compatible with gradient-based training than a hard rule that selects exactly one memory slot by an argmax operation?
Give both an optimization-based and a representational answer.
\end{exercise}

\subsection*{Notes and references}

The attention mechanism entered modern deep learning through sequence-to-sequence learning, where it was introduced as a remedy for the fixed-vector bottleneck of encoder--decoder recurrence.
Work by Bahdanau, Cho, and Bengio was especially influential in making this point explicit for neural machine translation.
Luong and collaborators developed closely related attention formulations, and the language of alignments helped connect neural translation to earlier ideas from statistical machine translation.
From a broader perspective, attention can be understood as differentiable content-based retrieval over an external set of representations.
That viewpoint became central once later architectures were built almost entirely from attention, but its pedagogical meaning is already clear here: attention solved a precise memory-access problem before it became the organizing principle of transformers.
